\noindent\textbf{Reading a translated passage and writing its translation place different demands on a visual language model. We held the bilingual glyph correspondences and neural writer fixed, and learned how the reader moves between source positions while emitting a translation. Four passes over 497,775 pictured sentence pairs in each direction took 42 seconds in total on one RTX 4090 D. On four previously unused books, both the ordered and unordered models retrieved the correct original passage for 91 of 96 images among 48 candidates. In free translation, all three visual writers stopped on every request. None earned a complete-meaning credit on those 96 passages. A quantized text model earned 81 credits under the same automated assessment.}\par
\begin{multicols}{2}
\section*{Learning where to read next}
A bilingual lexicon recognizes related passages even when its source strips are permuted. We introduced a latent position that can stay, move back one strip, move forward one or two, or return to a uniformly sampled position. A forward--backward fit estimated these probabilities from aligned books and parallel sentences. Emissions use the same fixed printed glyph repertoire in every condition. Free writing sums over source positions and possible boundaries between target strips. This construction adapts a classical alignment model to glyph prediction.\cite{vogel}
\end{multicols}
\displayfigure{Ordered_glyph_alignment.pdf}{What changes when a writer follows source order}{Top, the ordered writer adds the aspect marker \emph{le}, making its Chinese answer match the reference exactly. Both answers express the finding event. Below, the ordered writer recovers \emph{cold} and loses the speaker's name. The panels show the first newly exact controlled answer and the first changed English-to-Chinese book answer in test order. Each strip is an original source or output image.}
\begin{multicols}{2}
\section*{Control over a pictured sentence}
On 236 English sentence pairs that differed in negation, the neural, unordered and ordered writers produced both Chinese answers exactly 195, 207 and 205 times. On 118 sets that also exchanged the participants, they scored 84, 92 and 91. On the earlier 48-way book search, source movement raised first-place retrieval from 93/96 to 94/96.
\section*{Composition on a new book}
The ordered writer reproduced 236 of 1,416 distinct script-bearing reference strips, summed over requests. In an exploratory equal-length comparison, 18 of its 96 answers repeated a three-strip sequence, compared with one reference. An automated judge credited 63 supplied translations and no unrelated foil. These comparisons cover two printed languages and four works. \repo{references/ordered_alignment_result.md}{Methods, raw images and assessments}.
\end{multicols}
